\section{Introduction and the sharp observation theorem}
\label{sec:v83-main}
An action occurs in a clocked observation as the zero of an affine
factor $T-W$. A selective detector multiplies that factor by an unknown
function of the deadline. We ask which detector spaces permit recovery
of the zero after both total mass and component identity have been lost.
The central object is the logarithmic action curve modulo the detector
space. Its secants encode two unknown actions; its tangents control
regular inversion. Positivity of a contrast is useful, but is not a
necessary feature of an identifiable observation model.

We first give the sharp polynomial theorem. We then pass from a sampled
annihilator to a criterion on the detector \emph{function space} itself.
Theorem~\ref{thm:v84-design} turns regular separation of its secants into
a finite stable deadline design. Theorem~\ref{thm:v84-generic} gives
a uniform $d+5$ generic-clock bound for analytic spaces.
Theorem~\ref{thm:v84-poles} supplies an
explicit bound for prescribed logarithmic singularities, and
Proposition~\ref{prop:v84-no-positive} gives identification when no
positive annihilator exists at any set of clocks. The complete fibres
at one clock below the polynomial threshold are described in
Theorem~\ref{thm:v84-fibre}. The finite-record covering theorem and the
geometric applications are separate consequences of the resulting
pointwise inverses.

\subsection{Observation model}
Let $T_1<\cdots<T_J$ be real deadlines and let $Q$ be a compact connected
smooth domain, possibly with boundary or corners, with a fixed finite
coordinate atlas. Fix an integer $q\ge0$. At $z\in Q$ the observations
are positive projective classes of $u\times v$ matrices
\begin{equation}\label{eq:v83-model}
 M_j=U\diag\left\{a_b\exp\left(\sum_{\nu=1}^q
 \theta_{b,\nu}T_j^\nu\right)(T_j-W_b)\right\}_{b=1}^B V^{\mathsf T},
 \qquad a_b>0,\quad W_b<T_1.
\end{equation}
The columns of $U,V$ are probability vectors and have full column rank.
The channels may depend on $z$ but not on $j$. The two readings are
independent conditional on the component and endpoint. The visible count
$B\le\min(u,v)$ is not supplied. The normalized matrices are
\begin{equation}\label{eq:v83-probability}
 P_j=\frac{M_j}{\one^{\mathsf T}M_j\one}
     =U\diag(\pi_{j,b})V^{\mathsf T}.
\end{equation}
The joint matrix, not each row or column, is normalized. All coefficients
are locally $C^m$, $m\ge0$; local enumerations may differ by permutations.
The tuple of a component is
$(W_b,U_{:b},V_{:b},\log a_b,\theta_{b,1},\ldots,\theta_{b,q})$.
It defines a finite covering of $Q$, not necessarily a connected one.
No enumeration or trivialization is observed. Put
\begin{equation}\label{eq:v83-center}
 \alpha_b=\log a_b-B^{-1}\sum_c\log a_c,\qquad
 \vartheta_{b,\nu}=\theta_{b,\nu}-B^{-1}\sum_c\theta_{c,\nu},
 \qquad \Xi=(W,U,V,\alpha,\vartheta).
\end{equation}
These are global $C^m$ fields on the covering: the fibrewise sums are
permutation invariant and the local centred fields have the same
transition permutations as the original tuples.

\begin{theorem}[Sharp selective-detector rigidity]\label{thm:v83-main}
Suppose \eqref{eq:v83-model} holds. At every endpoint assume the tuples
$(W_b,\theta_{b,1},\ldots,\theta_{b,q})$ are pairwise distinct and the
actions are not all equal. Any $q+3$ distinct deadlines determine $B$,
the component covering and $\Xi$, up to a covering isomorphism over $Q$.
The complete continuous ambiguity is
\begin{equation}\label{eq:v83-gauge}
 a_b\mapsto h(z)a_b,\qquad
 \theta_{b,\nu}\mapsto\theta_{b,\nu}+\rho_\nu(z),
\end{equation}
where $h>0$ and $h,\rho_\nu$ are $C^m$ functions on $Q$, common to all
components of a fibre. No action changes. On simply connected $Q$ the
covering is trivial and two global enumerations differ by one constant
permutation.

On each reference neighbourhood specified in
Theorem~\ref{thm:v83-global}, the inverse has a locally Lipschitz
$C^m$ extension to real matrix fields, in the ambient centred affine
parameter space modulo the same covering isomorphisms. There is no
loss of spatial derivatives. The rank and covering stratum of this
extension are fixed; Theorem~\ref{thm:v84-mesh} separately treats their
recovery from finite records under quantitative signal assumptions.

The count is optimal over this observation class. With $q+2$ deadlines,
strictly positive two-component full-rank models with unequal actions
have an action-changing continuous ambiguity. In particular, four
clocks are necessary and sufficient in the worst log-affine case.
\end{theorem}
The optimality is a worst-case statement, not a requirement for every
larger model. Invisible duplicates and arbitrarily weak components
cannot be counted by a fixed signal threshold. Equal actions among
some components are allowed when other coefficients separate them and
one unequal-action pair remains.

\subsection{Relation to earlier methods}
The latent spectral step belongs to the established theory of
multiway identification and observable operators
\cite{V83AMR,V83AHK,V83BJR}. Here the deadline is a controlled index,
the two channels are unknown and may vary with the endpoint, and the
weights follow a restricted action--detector law. We use exact
full-rank and joint-separation hypotheses rather than a generic
identifiability assertion. The matrix factorization itself is not
claimed as a new tensor theorem.

After that step, the scalar data have the form
\begin{equation}\label{eq:v83-scalar}
 f(T)=P(T)+\log(T-x)-\log(T-y),\qquad P\in\cP_q,
 \quad x,y<T_1.
\end{equation}
We write $L_x(T)=\log(T-x)$ throughout.
Divided differences and Chebyshev systems are classical
\cite{V83deBoor,V83KS}. They give the two-moment proof below, but not
the positive-kernel-free example of Section~\ref{sec:v84-spaces}.
Our general criterion is a regular secant-separation criterion, not a
new characterization of arbitrary moment measures. The finite-design
argument uses compactness; the explicit logarithmic-family bound uses
rational zero counting and residues. Finite-rate-of-innovation and
Prony methods also recover finite singular supports
\cite{V84VMB,V84BY}; their linear moment or convolution samples are not
our projective matrices or sampled logarithms with an additive unknown
function space. We make no claim to have replaced those general theories.

The covering construction uses local spectral perturbation and covering
space theory \cite{V84Kato,V84Hatcher}. Its additional task is to match
actual channel columns across independently chosen nonorthogonal minor
coordinates. The finite-record theorem supplies edge permutations and
monodromy from observations; it does not assume a reference covering.
Generating functions and exact-lift composition are standard
\cite{V83MW,V83Geiges}. Section~\ref{sec:v84-observed-geometry} replaces
the prior common-lift premise by literal identities of the reconstructed
lifts on an observed invariant domain. For the surface application,
Pestov--Uhlmann's boundary-distance rigidity theorem supplies the metric
step \cite{V84PU}; our conclusion is the reduction from uncalibrated
selective records to its distance data. The local and global partial-data
metric results in \cite{V84SUV} solve different geometric inverse
problems and are not inferred from our finite-dimensional clock theorem.

\section{Oriented intervals and detector quotients}
\label{sec:v83-interval}
\begin{lemma}[Two moments of an interval]\label{lem:v83-interval}
Let $I$ be an open interval, $w:I\to(0,\infty)$ continuous, and
$r:I\to\R$ continuous. On $x\ne y$ the map
\[
 \mathcal M(x,y)=\left(\int_y^xw(s)\dd s,
                         \int_y^xw(s)r(s)\dd s\right)
\]
is injective on all oriented intervals if and only if $r$ is strictly
monotone. If $w,r$ are smooth and $r$ is strictly monotone, it is a local
smooth diffeomorphism and
\begin{equation}\label{eq:v83-moment-jacobian}
 \det D\mathcal M(x,y)=w(x)w(y)\bigl(r(x)-r(y)\bigr).
\end{equation}
\end{lemma}
\begin{proof}
Choose $F'=w$, put $u=F(y)$, $d=F(x)-F(y)$, and $R=r\circ F^{-1}$
on $I'=F(I)$. The first moment is $d$. For $d>0$ the second is
$H_d(u)=\int_u^{u+d}R(t)\dd t$. If $R$ is strictly increasing,
$H'_d(u)=R(u+d)-R(u)>0$; strict decrease reverses the inequality.
The two moments determine $d,u$, hence $x,y$. The sign of $d$ separates
orientations, and reversing endpoints handles $d<0$.

Conversely, injectivity makes each $H_d$ injective on the interval
$\{u:u,u+d\in I'\}$, hence strictly monotone. A constant interval of
$R$ would give a constant interval of $H_d$ for small $d$, and is
excluded. Choose $u_0<u_1$ with $R(u_0)\ne R(u_1)$. Exhaust $I'$ by
compact intervals whose interiors contain these two points. On any
such compact interval, $H_d/d\to R$ uniformly as $d\downarrow0$.
For all sufficiently small $d$ the sign of
$H_d(u_1)-H_d(u_0)$ is thus fixed by $R(u_1)-R(u_0)$, and fixes the
monotonicity orientation of $H_d$ on its entire interval of definition.
For arbitrary $v_0<v_1$ in $I'$, choose an exhaustion interval containing
all four points in its interior. Passing to the limit in the inequality
for $H_d(v_0),H_d(v_1)$ shows that $R$ has the same nonstrict
monotonicity on all of $I'$. Equality at two distinct points would make
it constant between them, already excluded. Thus it is strictly
monotone. Differentiation gives the determinant.
\end{proof}
This is a theorem about moving endpoints and a known density, not
about recovery of arbitrary measures from two moments. Its global
conclusion does not follow merely from its nonsingular Jacobian.

\subsection{A sampled positive-kernel criterion}
Let $J=d+2$, let $\cD\subset\R^J$ have dimension $d$ and contain
$\one$, and choose a basis $c^0,c^1$ of $\cD^\perp$. Define
\begin{equation}\label{eq:v83-kernels}
 k_a(s)=\sum_{j=1}^J\frac{c_j^a}{T_j-s},\qquad a=0,1.
\end{equation}
Assume the kernel space has a member positive on
$I\subset(-\infty,T_1)$, and choose $k_0>0$.
\begin{theorem}[Positive-kernel detector classification]
\label{thm:v83-classification}
For $\ell_j=e_j+\log(T_j-x)-\log(T_j-y)$, $e\in\cD$ and
$x\ne y$ in $I$, the ordered actions and sampled nuisance vector are
globally identifiable if and only if $k_1/k_0$ is strictly monotone.
Equivalently,
\begin{equation}\label{eq:v83-chebyshev}
 k_0(s)k_1(t)-k_1(s)k_0(t)
\end{equation}
has one strict sign for $s<t$ in $I$. The condition is independent of
the annihilator basis. Its failure gives distinct nondegenerate action
intervals and nuisance vectors with exactly the same scalar data.
\end{theorem}
\begin{proof}
Since $\log(T-x)-\log(T-y)=-\int_y^x(T-s)^{-1}\dd s$,
the contrasts $(-c^0\cdot\ell,-c^1\cdot\ell)$ are the two interval
moments. Equal contrasts are sufficient as well as necessary for the
logarithmic action vectors to differ by an element of $\cD$.
Apply Lemma~\ref{lem:v83-interval}, then subtract to recover $e$.
A basis change multiplies \eqref{eq:v83-chebyshev} by its nonzero
determinant. Between two bases with positive first kernel, the new
ratio is a real fractional-linear transform of the old one with a
nowhere-zero denominator; its derivative with respect to that ratio
has constant nonzero sign. This also proves ratio-level invariance.
\end{proof}
The hypothesis on a positive kernel is not universal. General function
spaces and higher-dimensional annihilators are treated below. This
scalar theorem alone supplies neither spectral separation of a latent
mixture nor recovery of function coefficients from a rank-deficient
evaluation map.

\subsection{Sharp polynomial clock rigidity}
For distinct nodes use
\[
 [t_0,\ldots,t_n]f=\sum_i\frac{f(t_i)}{\prod_{h\ne i}(t_i-t_h)}.
\]
It annihilates polynomials of degree less than $n$, and induction from
the first difference gives
\begin{equation}\label{eq:v83-cauchy-dd}
 [t_0,\ldots,t_n]\frac1{T-s}=\frac{(-1)^n}{\prod_i(t_i-s)}.
\end{equation}
\begin{theorem}[The sharp scalar inverse]\label{thm:v83-scalar}
For $q\ge0$, $q+3$ distinct values of \eqref{eq:v83-scalar} determine
$x\ne y<T_1$ and $P$ globally, with a locally smooth inverse. With
only $q+2$ values, action-changing one-parameter families give the same
data. Their entire fibres are parametrized in
Theorem~\ref{thm:v84-fibre}.
\end{theorem}
\begin{proof}
Take the consecutive order-$q+1$ differences
\begin{align}
 d_0&=(-1)^{q+2}[T_1,\ldots,T_{q+2}]f,
 &w_0(s)&=\prod_{i=1}^{q+2}(T_i-s)^{-1},\label{eq:v83-contrasts}\\
 d_1&=(-1)^{q+2}[T_2,\ldots,T_{q+3}]f,
 &w_1(s)&=\prod_{i=2}^{q+3}(T_i-s)^{-1}.\nonumber
\end{align}
They equal $\int_y^xw_0,\int_y^xw_1$. These two annihilators are
independent and span the polynomial evaluation annihilator. Moreover
\begin{equation}\label{eq:v83-ratio}
 \frac{w_1(s)}{w_0(s)}=\frac{T_1-s}{T_{q+3}-s},\qquad
 \left(\frac{w_1}{w_0}\right)'(s)
 =\frac{T_1-T_{q+3}}{(T_{q+3}-s)^2}<0.
\end{equation}
The interval lemma recovers $x,y$; subtraction and interpolation recover
$P$. It also proves smoothness. Directly, the scalar Jacobian has
columns $1,T,\ldots,T^q,-(T-x)^{-1},(T-y)^{-1}$. Multiplying a row
relation by $(T-x)(T-y)$ gives a polynomial of degree at most $q+2$
vanishing at $q+3$ nodes. Evaluation at $x,y$ forces the two pole
coefficients to vanish, and then the polynomial coefficients vanish.
At $q+2$ nodes the columns $1,T,\ldots,T^q,-(T-x)^{-1}$ alone are
independent by the same argument with denominator $T-x$. The implicit
function theorem solves for these parameters while a nearby $y$ varies.
Strict probability and action margins persist on a small segment.
\end{proof}
For $d_0>0$, choose $F'=w_0$ and solve the strictly decreasing equation
\begin{equation}\label{eq:v83-monotone-solve}
 \int_u^{u+d_0}\frac{w_1(F^{-1}(t))}{w_0(F^{-1}(t))}\dd t=d_1.
\end{equation}
Then $y=F^{-1}(u)$ and $x=F^{-1}(u+d_0)$. The condition
$u,u+d_0\in F((-
\infty,T_1))$ gives the bracketing interval. Reverse orientation for
negative $d_0$. No action sorting is used. Once $y$ is known, the
single contrast $\int_y^xw_0$ is strictly increasing in $x$, including
$x=y$, so $q+2$ nodes identify every remaining component. The first
contrast vanishes exactly at equal actions.
